\subsection{结果与机制解释}
表 \ref{tab:q1_result} 给出了 2025 年的总量标定结果与局地对照结果。四大家鱼总量与题目给出的恢复信息接近；$\mathrm{CPUE}^*$ 虽低于 2.2 的湖北段对照值，但恢复方向一致。
\begin{table}[htbp]
\centering
\caption{问题一的 2025 年参数标定结果与局地对照}
\label{tab:q1_result}
\begin{tabularx}{0.95\textwidth}{>{\raggedright\arraybackslash}p{3.0cm} >{\raggedright\arraybackslash}p{2.6cm} >{\raggedright\arraybackslash}p{2.6cm} >{\raggedright\arraybackslash}p{2.8cm}X}
\toprule
指标 & 观测值 & 模拟值 & 相对误差 & 说明 \\
\midrule
四大家鱼总量倍率 & 1.800 & 1.805 & 0.26\% & 对齐 \\
$\mathrm{CPUE}^*$ 倍率 & 2.200 & 1.838 & 16.47\% & 一致 \\
\bottomrule
\end{tabularx}
\end{table}

两项指标的方向相同，幅度却不同。四大家鱼总生物量基本复现用于标定的1.8倍，表明单区模型在全流域平均尺度上与标定目标一致；$\mathrm{CPUE}^*$低于湖北江段的2.2倍，局地渔获口径可能带来另一种增幅。两者无法用同一误差解释。

禁渔情景下，四大家鱼总量与$\mathrm{CPUE}^*$上升，基础资源末值为初值的0.204倍，水草最低值出现在第10年；食压指数同时增至4.213。这些读数表明，鱼群恢复与底层资源改善在模型中未同步发生。0.204和$\Phi$均仅为归一化单区指标，不代表全流域实际生物量或能量金字塔的绝对比例。$\Phi$持续上升与底层资源下降共同表明，消费者增长快于资源供给；捕捞压力解除后，中层鱼类的摄食压力和种间竞争随之增强，呈现"鱼类资源恢复与基础资源承压并存"。

\begin{figure}[htbp]
\centering
\includegraphics[width=0.92\textwidth]{fig01_q1_foodweb.png}
\caption{问题一：标定后的禁渔与对照情景及观测约束}
\label{fig:q1}
\end{figure}

